% 01_overview.tex — Section 1: System Overview & Architecture
% Phase 2 content migration — Task #178, Cycle 54 NVB-1, 2026-03-20
% Sources: hardware/sipm_fpga/ARCHITECTURE_SUMMARY.tex,
%          hardware/hardware_readiness_report.tex
% Law 16: quantitative claims carry source locator comments.
% Status: POPULATED (Phase 2)

\section{System Overview \& Architecture}
\label{sec:overview}

The NV bolometer is a cryogenic particle detector that uses nitrogen-vacancy (NV) centers
in diamond to read out the Meissner-field collapse of a superconducting NbN thin-film pixel.
When a particle deposits energy exceeding the thermal threshold $E_\mathrm{th}$ in the NbN
film, the film transitions from the superconducting (Meissner-screened) to the normal state.
The change in local magnetic field at the NV layer is detected optically via a shift in the
optically detected magnetic resonance (ODMR) frequency of the NV ensemble~\cite{Doherty2013}.
% Doherty2013 p.1--4 (NV centre review: ODMR principle, ZPL at 637 nm)

\subsection{Sensing Principle}

The NV$^-$ ground state is a spin-1 triplet with zero-field splitting $D \approx 2877.5$\,MHz
at $T = 12$\,K (This work; CryoLab measurement).
% internal: hardware_readiness_report.tex Sec.~2 (D = 2877.5 MHz at 12 K confirmed in lab)
In the presence of a static bias field $B_\mathrm{bias}$, the $m_s = \pm 1$ sub-levels split
by $2\gamma_e B_\mathrm{bias}$, where $\gamma_e / 2\pi = 28.0$\,GHz/T~\cite{Doherty2013}.
% Doherty2013 p.~33, Eq.~(A3) (gyromagnetic ratio of NV electron spin)
When the NbN film is superconducting, Meissner screening perturbs the local field at the NV
layer. When the film transitions to the normal state (either thermally or via a particle
absorption event), this perturbation collapses. The resulting change in field $\Delta B_z$
produces an ODMR frequency shift $\Delta\nu = \gamma_e \Delta B_z / \pi$, which is the
detection signal (OQ-1 mechanism; see Section~\ref{sec:oq1}).

\subsection{Milestone Architecture}

\begin{table}[ht]
\centering
\caption{NV bolometer milestone roadmap. Energy thresholds are from the threshold-detector
  model (Section~\ref{sec:oq5}). \% internal: oq5\_energy\_resolution.tex Eq.~(2).
  D0 uses the existing Montana cryostat (11\,K base, This work).}
\label{tab:milestones}
\small
\begin{tabular}{@{}p{1.4cm}p{2.5cm}p{1.8cm}p{2.5cm}p{4.5cm}@{}}
\toprule
Stage & SC Film & $T_\mathrm{op}$ & Readout & Primary Goal \\
\midrule
D0-Easy & NbN 200\,nm / sapphire & 12\,K & CW-ODMR & Meissner--Zeeman proof of concept \\
D0-Opt  & NbN, $1\,\mu\text{m}^2$ pixel & 12\,K & CW-ODMR & Energy-resolving demonstration \\
D1      & NbN/Al $+$ Ge absorber & 4.2\,K & SiPM + FPGA & Single-pixel bolometric detection \\
P1      & Al 20\,nm / graphite & 0.75\,K & SiPM + FPGA & Sub-eV dark matter detection \\
\bottomrule
\end{tabular}
\end{table}

\subsection{Full D1 Readout Chain}

The selected D1 readout architecture (Option~A, \texttt{ARCHITECTURE\_SUMMARY.tex}) is:
% internal: hardware/sipm_fpga/ARCHITECTURE_SUMMARY.tex (architecture decision 2026-03-18)
\begin{equation*}
  \underbrace{\text{S14420 SiPM}}_{\text{photon detection}}
  \;\to\;
  \underbrace{\text{MAX40658 TIA}}_{\text{current--voltage}}
  \;\to\;
  \underbrace{\text{ADCMP601 comp.}}_{\text{pulse discrimination}}
  \;\to\;
  \underbrace{\text{Nexys A7 FPGA}}_{\text{timing \& control}}
  \;\to\; \text{PC}.
\end{equation*}
The FPGA runs pulse sequencing, ADF4351 SPI frequency control,
carry-chain TDC photon timestamping at 17.5\,ps resolution~\cite{Lim2025},
% Lim2025 Table 4 p.13 (4-tap SCSC LSB = 17.530 ps, Artix-7 CARRY4 chain)
histogram accumulation, and PC UART communication on a single board.

\subsection{Current Laboratory State}

\begin{itemize}
  \item \textbf{NV diamond} (HPHT ensemble): $D = 2877.5$\,MHz at 12\,K; ODMR contrast
    6--7.5\%; $T_1 \approx 230$\,ms at 20--65\,K;
    sensitivity $\sim\!14\,\mu\text{T}/\sqrt{\text{Hz}}$ at 12\,K.
    (This work; CryoLab ODMR campaign.)
  \item \textbf{Montana cryostat}: 11\,K base; ODMR data collected at 12--250\,K.
  \item \textbf{ADF4351 synthesizer}: 35\,MHz--4.4\,GHz, $+5$\,dBm, SPI-controlled~\cite{ADF4351datasheet}.
    % ADF4351datasheet p.4 Table 1 (output power, frequency range)
  \item \textbf{S14420 SiPM}: in hand; cryogenic characterization pending
    (Section~\ref{sec:sipm-cryo}).
  \item \textbf{NbN thin film}: deposition pending at UT Arlington Shimadzu NRC
    (Section~\ref{sec:nbn-sourcing}).
\end{itemize}
